\chapter{SpriteBatch and 2D Rendering}
\label{ch:spritebatch}

\cnaclass{SpriteBatch} is the API most CNA code touches most often, and --- because it is
also the one class every one of the eleven backends in Part~IV must implement --- the class
whose bugs and fixes best illustrate what ``backend parity'' actually costs to achieve. This
chapter covers its API surface, its supporting types (\cnaclass{SpriteSortMode},
\cnaclass{SpriteEffects}, \cnaclass{SpriteFont}), and the specific, real bugs found while
verifying it, because the bugs are as instructive as the API itself.

\section{Begin, Draw, End}

\cnaclass{SpriteBatch} inherits \cnaclass{GraphicsResource} and offers five \texttt{Begin()}
overloads of increasing specificity, each defaulting whatever the caller omits:
\texttt{Begin()} alone defaults to \texttt{SpriteSortMode::Deferred},
\cnaclass{BlendState::AlphaBlend}, a \texttt{LinearClamp} sampler, no depth-stencil state,
\texttt{RasterizerState::\allowbreak{}CullCounterClockwise}, and an identity transform; the fuller
overloads let a caller override the sort mode and blend state, the three optional state
objects (nullptr falls back to the same defaults), a custom \cnaclass{Effect} (nullptr falls
back to an internal default \cnaclass{SpriteEffect}), and finally an explicit transform
\cnaclass{Matrix} applied to every sprite in the batch before projection.

\texttt{Draw} itself has nine overloads: six are the canonical XNA forms ---
position-or-destination-rectangle, an optional source rectangle (represented as
\texttt{std::optional<Rectangle>} rather than C\#'s nullable \texttt{Rectangle?}), color, and
progressively more control over rotation, origin, scale (as either a uniform \texttt{float} or
a non-uniform \cnaclass{Vector2}), \cnaclass{SpriteEffects}, and layer depth --- and three are
\texttt{NOXNA} raw texture-and-rectangle convenience forms layered on top.
\texttt{DrawString} adds six more overloads (\texttt{string} or \texttt{StringBuilder}, each in
a simple 4-argument form, a uniform-scale 9-argument form, and a non-uniform-scale
(\cnaclass{Vector2} scale) 9-argument form).

\section{The Draw overload family, and a worked rotation example}

The nine \texttt{Draw} overloads split cleanly into two families by their second parameter ---
a \cnaclass{Vector2} \texttt{position} (four overloads, increasingly parameterized) or a
\cnaclass{Rectangle} \texttt{destinationRectangle} that scales the texture to fit (three
overloads) --- plus the two \texttt{NOXNA} raw \texttt{float x, float y} convenience forms
already noted. Every real (non-\texttt{NOXNA}) overload takes an explicit \texttt{Color} tint
--- there is no overload that omits it, unlike the \texttt{NOXNA} convenience forms. The table
below abbreviates each overload's parameter names for width (\texttt{tex} for \texttt{texture},
\texttt{pos} for \texttt{position}, and so on) --- the full, exact parameter names for every
overload are given in the fully-parameterized worked example immediately below the table.

\begin{center}
\begin{longtable}{>{\raggedright\arraybackslash}p{0.30\textwidth}>{\raggedright\arraybackslash}p{0.58\textwidth}}
\toprule
\textbf{Family member} & \textbf{Adds, relative to the simplest form in its family} \\
\midrule
\endhead
\texttt{Draw(tex, pos, color)} & The simplest real-XNA form: whole texture, no rotation, unit
  scale. \\[4pt]
\texttt{Draw(tex, pos, srcRect, color)} & An optional source sub-rectangle
  (\texttt{std::optional<Rectangle>}, not C\#'s nullable \texttt{Rectangle?}) --- the
  sprite-sheet-cell pattern. \\[4pt]
\texttt{Draw(tex, pos, srcRect, color, rot, origin, scale, fx, depth)} & Uniform
  \texttt{float} scale. \\[4pt]
\texttt{Draw(tex, pos, srcRect, color, rot, origin, scale, fx, depth)} & Same, but
  \texttt{scale} is a non-uniform \cnaclass{Vector2} --- these two fully-parameterized
  overloads differ \emph{only} in that one field's type. \\[4pt]
\texttt{Draw(tex, destRect, color)} & Whole texture, scaled to fill \texttt{destRect}. \\[4pt]
\texttt{Draw(tex, destRect, srcRect, color)} & Adds the source sub-rectangle. \\[4pt]
\texttt{Draw(tex, destRect, srcRect, color, rot, origin, fx, depth)} & Full control,
  destination-rectangle family --- note there is no separate \texttt{scale} parameter here at
  all, since the destination rectangle already implies it. \\
\bottomrule
\end{longtable}
\end{center}

\texttt{origin} is the parameter every one of the fully-parameterized overloads shares, and it
is worth a worked example precisely because \S\ref{sec:sdl-renderer-two-bugs} below names a
real bug caused by getting a pivot convention wrong: \texttt{origin} is a point \emph{in
texture space} (not screen space) that both the rotation pivot and the drawn position are
measured relative to. Rotating a $64\times64$ sprite around its own center, rather than its
top-left corner, means passing that texture's own center as \texttt{origin}:

\begin{lstlisting}
Vector2 center(texture.getWidthProperty() / 2.0f, texture.getHeightProperty() / 2.0f);
float angle = totalSeconds; // radians; grows every frame for a continuous spin

spriteBatch_->Draw(
    texture, spritePosition, std::nullopt, Color::White,
    angle, center, 1.0f, SpriteEffects::None, 0.0f);
\end{lstlisting}

Passing \texttt{Vector2::Zero} for \texttt{origin} instead --- the easy mistake, since it is
also the type's default-constructed value --- rotates the sprite around its top-left corner,
which reads as the sprite ``orbiting'' \texttt{spritePosition} in a circle rather than spinning
in place; the visual symptom is easy to misdiagnose as a rotation-math bug when the actual
cause is simply the wrong \texttt{origin}.

\section{SpriteSortMode and SpriteEffects}

\cnaclass{SpriteSortMode} has five values: \texttt{Deferred}, \texttt{Immediate},
\texttt{Texture}, \texttt{BackToFront}, \texttt{FrontToBack}. Only the latter two actually
honor a sprite's \texttt{layerDepth} argument --- the others ignore it entirely, by design,
matching XNA's own documented sort-mode semantics.

\cnaclass{SpriteEffects} is where this chapter's first concrete, still-open gap lives:
it is declared as a plain \texttt{enum class} --- \texttt{None=0}, \texttt{FlipHorizontally=1},
\texttt{FlipVertically=2} --- with \emph{no bitwise operators defined at all}. In real XNA,
\texttt{SpriteEffects} is a \texttt{[Flags]} enum, and combining both flags
(\texttt{FlipHorizontally | FlipVertically}, value 3) is a normal, supported call. In CNA,
reaching value 3 requires an explicit \texttt{static\_cast<SpriteEffects>(3)}, and --- more
tellingly --- \texttt{SpriteBatch::DrawString}'s own per-axis direction lookup tables (added
when the flip-mirroring bug below was fixed) deliberately only have three entries, indexed
0 through 2. Value 3 is functionally unreachable through the normal calling convention. This
is flagged directly in \texttt{docs/spritefont-support.md} \S7 as a genuine, minor,
still-open API-completeness gap, not a crash risk --- a combined horizontal-and-vertical
flip of drawn text simply is not expressible today.

\section{A real, shared bug: SpriteEffects flip only affected sampling, not layout}

Chapter~\ref{ch:design-philosophy} promised that this book would name specific, checkable
bugs rather than paraphrase a changelog, and this one is worth naming in full because it
affected every backend at once, which is unusual --- most of the bugs catalogued in Part~IV
are backend-specific. Before Task~694, a \texttt{SpriteEffects} flip on
\texttt{DrawString} only flipped each individual glyph's own texture sampling in place; it
never mirrored the \emph{sequence or position} of the glyphs themselves, which is what FNA's
real implementation does via a pair of axis-direction/axis-mirrored terms. The fix lives in
the shared, backend-agnostic \texttt{SpriteBatch.cpp} --- meaning it was one fix that
corrected the behavior identically on every one of CNA's graphics backends simultaneously,
rather than something Part~IV's chapters need to revisit backend by backend.

\section{SpriteFont}

\cnaclass{SpriteFont} is deliberately not a \cnaclass{GraphicsResource} --- a texture atlas,
glyph bounds, cropping data, a character set, line spacing, character spacing, kerning data,
and an optional default character. Its constructor is \texttt{NOXNA} and publicly callable,
a direct consequence of Chapter~\ref{ch:content-and-assets}'s content model: real XNA's
\cnaclass{SpriteFont} constructor is \texttt{internal}, reachable only from the compiled
\texttt{.xnb} content pipeline's own \texttt{SpriteFontReader}; CNA has no such reader, so the
constructor has to be public for anything to construct a font at all.

\texttt{MeasureString} exists for both \texttt{string} and \texttt{StringBuilder} --- the
\texttt{StringBuilder} overload was entirely missing until Task~421/423, and is now a
one-line forward to the \texttt{string} implementation, a simplification available in
C\texttt{++} precisely because there is no garbage-collection-pressure argument (the reason
FNA's own C\# implementation duplicates the algorithm instead of forwarding) against sharing
code here. A faithfully-preserved, slightly surprising detail: a trailing \texttt{"\textbackslash n"}
adds a full second empty line's worth of height to the measured result ---
\texttt{Y = 2 \(\times\) LineSpacing}, not \texttt{1 \(\times\)} --- because both the newline
handler and the loop's own unconditional final height-add fire. This matches real XNA's
behavior exactly and is preserved rather than ``fixed,'' consistent with
Chapter~\ref{ch:design-philosophy}'s fidelity-over-taste rule.

Character encoding deserves its own note: CNA's \texttt{String} is UTF-8
(\texttt{std::string}), and \cnaclass{SpriteFont} / \cnaclass{SpriteBatch} decode each glyph via
an internal \texttt{DecodeUtf8CodePoint} helper. Because \texttt{charcs} (the glyph-key type,
per Chapter~\ref{ch:design-philosophy}'s type-alias table) is 16 bits wide, only Basic
Multilingual Plane code points can serve as glyph keys at all; an invalid or truncated UTF-8
sequence decodes to \texttt{'?'} rather than throwing, and the byte index always advances by
at least one, so a malformed string can never cause an infinite decode loop.

A worked example ties \texttt{MeasureString} directly to the simplest \texttt{DrawString}
overload for the single most common real use --- centering a line of HUD text horizontally
around a screen-space point, which needs the string's own measured width before it can compute
where to start drawing:

\begin{lstlisting}
const std::string text = "Game Over";
Vector2 size = font.MeasureString(text);           // Vector2(width, height) in pixels
Vector2 centered(screenCenterX - size.X / 2.0f, screenCenterY - size.Y / 2.0f);

spriteBatch_->DrawString(font, text, centered, Color::White);
\end{lstlisting}

\section{What SpriteFont's support matrix actually says}

\texttt{docs/spritefont-support.md} verifies the property surface, \texttt{MeasureString},
single- and multi-glyph placement, newline handling, default-character fallback, single-axis
flip, and rotation/scale as pixel-correct on both \textbf{SDL\_Renderer} and \textbf{EasyGL}.
On \textbf{Vulkan} and \textbf{BGFX}, the underlying logic is the same shared C\texttt{++} ---
very likely correct by construction --- but has not been independently re-verified
specifically for \cnaclass{SpriteFont} (flagged again in
\texttt{docs/\allowbreak{}graphics-\allowbreak{}backend-\allowbreak{}feature-\allowbreak{}matrix.md},
Task~861). Combined-axis flip is not
representable on any backend at all, for the \cnaclass{SpriteEffects} reason above.

\section{Two real, fixed SDL\_RENDERER bugs, as a preview of Part~IV's method}
\label{sec:sdl-renderer-two-bugs}

Two of the fifteen real bugs found during the SDL\_RENDERER backend's own Phase~70 audit
(Chapter~\ref{ch:sdl-renderer} covers the full audit) live directly in \cnaclass{SpriteBatch}
rendering and are worth previewing here as a template for how this book reports a bug
throughout Part~IV: name the symptom, the root cause, and the fix. First,
\texttt{SDL\_RenderTextureRotated}'s pivot convention differed from XNA's own rotation-origin
convention; the fix offsets the destination rectangle rather than changing the rotation math
itself. Second --- more serious --- the \texttt{transformMatrix} argument to
\texttt{Begin()} was silently ignored \emph{entirely} on this backend; any code passing a
non-identity transform matrix (a common pattern for camera-style 2D scrolling) simply drew as
if it had passed identity. The fix added a new rendering path using
\texttt{SDL\_RenderTextureAffine}, used only when the transform actually differs from
identity, so the common (and cheaper) untransformed path is unaffected.

\section{Where this leaves a porting engineer}

\cnaclass{SpriteBatch} is, by a wide margin, the most exhaustively cross-backend-verified
class in CNA --- every backend chapter in Part~IV returns to it, because every backend has
to implement it, and because its bugs (an ignored transform matrix, a mis-scoped flip, a
silently-downgraded sampler filter) are exactly the class of bug that only surfaces under
real pixel-readback testing rather than ``it compiles and something appears on screen.'' If
you are porting existing 2D XNA/FNA game code (this book's migration chapter), this is also
the API you can trust most: it is the one every backend is held to the same standard on.
